\section{Sharp cyclicity}
\label{sec:curvature-graph}

We count fixed points by following the zeros of $P''$, which we call the
curvature of the return. Where both passages are active and the state is
large compared with the separation (the outer region), these zeros lie on a
separation graph with at most one positive critical point; estimates near the
thresholds (the threshold layer) fix the sign of $P''$ there. Together they
give the upper bounds on the whole interval $I$. We then construct
configurations attaining the bounds under~\eqref{eq:rank}.

An active three-halves term has curvature proportional to \(q^{-1/2}\).
In the outer region, where both passages are active, \(q=r^2\) and \(J=rP''\) remove
this singular factor, while \(v=\delta/q\) measures separation relative
to the state. With the aligned multiplier \(k(p)=1+b>0\), their first-order balance is
\[
 J=2c_0r+\frac{3k}{4}m-\frac{3kA}{8}v
       +O\bigl((r+|m|+|v|)^2\bigr).
\]
The strength $A\ne0$ makes the curvature equation transverse in $v$.
Exact division then gives a separation graph $\Delta$ with
$\Delta_r=r\mathcal W$; the closing curvature $c_0\ne0$ makes
$\mathcal W$ strictly monotone. These two reductions give at most one
positive critical point.

\subsection{The exact curvature-zero graph}

Let \(g_\nu\) be the signed-root companion in
Proposition~\ref{prop:quadratic-passage}. Its dependence on
\((p,m,\delta)\) is suppressed in the notation. Define smooth functions
\begin{equation}\label{eq:root-profiles}
 T_\nu(u)=g_\nu(u)/u,\quad
 \ell_\nu(u)=T_\nu(u)^2+\tfrac{u}{2}(T_\nu^2)'(u),
 \quad H_\nu(u)=\tfrac12\ell_\nu'(u).
\end{equation}
The quotient is filled in smoothly at \(u=0\), with \(T_\nu(0)=1\).
For \(u>0\),
\begin{equation}\label{eq:passage-derivatives}
 \Gamma_\nu'(u^2)=\ell_\nu(u),\qquad
 \Gamma_\nu''(u^2)=u^{-1}H_\nu(u),\qquad
 \ell_\nu(0)=1,\quad H_\nu(0)=\tfrac34 A_\nu.
\end{equation}
These are exact identities obtained by differentiating
\(\Gamma_\nu(u^2)=g_\nu(u)^2\).

For \(r>0\), put
\begin{equation}\label{eq:J-definition}
 q=r^2,\quad \delta=r^2v,\quad
 J(r,m,v,p)=rP_{m,r^2v}''(r^2,p).
\end{equation}
Choose a fixed small \(0<\eta\le v_0=1/4\), within the box of Lemma~\ref{lem:exact-value-division}. All profiles in the following
calculation are evaluated at \(\delta=r^2v\).
Write
\[
 U=T_1(r)^2,\qquad W=U-v,\qquad s_2=r\sqrt{W}.
\]
After a fixed common shrinking, \(W>1/2\) for \(|v|\le\eta\).
The two excursions are then both active and the exact downstream
penetration is \(y=r^2W>0\). Chain differentiation gives
\begin{align}
 Q&=r^2\{v+WT_2(s_2)^2\},&
 Q'&=\ell_2(s_2)\ell_1(r),\label{eq:Q-outer}\\
 Q''&=r^{-1}\{\ell_2(s_2)H_1(r)
       +W^{-1/2}H_2(s_2)\ell_1(r)^2\}.\label{eq:Q-curvature}
\end{align}
Consequently
\begin{align}\label{eq:J-exact}
 J={}&rR''(Q)[\ell_2(s_2)\ell_1(r)]^2\\
 &+R'(Q)\{\ell_2(s_2)H_1(r)
               +W^{-1/2}H_2(s_2)\ell_1(r)^2\}.\notag
\end{align}
Every expression on the right is a root profile, jointly smooth through
\(r=0\).

\begin{lemma}[Root curvature and its zero graph]\label{lem:curvature-profile}
The function \(J\) has a joint smooth extension on a common box
in \((r,m,v,p)\). At its root boundary,
\begin{equation}\label{eq:J-boundary}
 J(0,m,v,p)
 =\tfrac34 R'_{p,m,0}(0)
     \{A_1(p,m,0)+A_2(p,m,0)(1-v)^{-1/2}\}.
\end{equation}
In particular,
\begin{equation}\label{eq:J-transverse}
 J_v(0,0,0,p)=-\tfrac38 R'_{p,0,0}(0)A(p)\ne0.
\end{equation}
There is exactly one smooth graph \(v=V(r,m,p)\), with
\(|V|<\eta\), on which \(J=0\) in this box.
\end{lemma}

\begin{proof}
Formula~\eqref{eq:J-exact} proves smoothness, and
\eqref{eq:passage-derivatives} at \(r=s_2=0\) gives
\eqref{eq:J-boundary}. Differentiation in \(v\) gives
\eqref{eq:J-transverse}. First choose \(\eta\) so that \(J_v\) has
the strict sign \(-A\) and a uniform nonzero margin on the root
boundary. The signs of \(J(0,0,-\eta,p)\) and
\(J(0,0,\eta,p)\) are opposite. Compactness in the remaining
parameters permits one further shrinking of \(r,m\) on which these
properties persist throughout the box. The implicit function theorem,
existence by the endpoint signs, and strict monotonicity in \(v\)
give the asserted unique graph.
\end{proof}

The boundary value of the graph is exact:
\begin{equation}\label{eq:V-boundary}
 V(0,m,p)=1-\left(\frac{-A_2(p,m,0)}{A_1(p,m,0)}\right)^2,
 \qquad V_m(0,0,p)=\frac{2}{A(p)}.
\end{equation}
Indeed the ratio equals \(1-m/A_1(p,m,0)\); both individual strengths
retain their dependence on \(m\).

To bound the positive critical points of the separation graph
\(r\mapsto r^2V(r,m,p)\) on a fixed interval, divide first along
the cancelling slice and then in the independent strength variable.
Finite cancellation gives
\begin{equation}\label{eq:Q-aligned}
 Q_{0,0}(q,p)=q+q_+^{5/2}L(\sqrt{q_+},p).
\end{equation}
For positive \(q=r^2\) it follows, by differentiation in \(q\),
that
\begin{equation}\label{eq:center-curvature}
 P_{0,0}''(r^2,p)=2c_0(p)+rN(r,p),\qquad
 c_0(p)=\tfrac12R_{p,0,0}''(0),
\end{equation}
where \(N\) is jointly smooth. The offset and linear jets are allowed
to vary in \(p\) throughout this identity.

Let \(V_0(r,p)=V(r,0,p)\). Since \(J(r,0,V_0,p)=0\), the mean-value
identity in the independent variable \(v\) yields
\[
 0=r\{2c_0+rN\}
       +V_0\int_0^1J_v(r,0,tV_0,p)\,dt.
\]
The integral has a nonzero common margin. Therefore \(V_0=rB_\Delta(r,p)\)
with \(B_\Delta\) smooth and
\begin{equation}\label{eq:B-boundary}
 B_\Delta(0,p)=\frac{16c_0(p)}{3R'_{p,0,0}(0)A(p)}.
\end{equation}
Hadamard division in \(m\) now gives the exact identity
\begin{equation}\label{eq:V-division}
 V(r,m,p)=mH_\Delta(r,m,p)+rB_\Delta(r,p),\qquad
 H_\Delta(0,0,p)=\frac{2}{A(p)}.
\end{equation}

Define the separation graph
\begin{equation}\label{eq:Delta-definition}
 \Delta(r,m,p)=r^2V(r,m,p).
\end{equation}
Its derivative factors exactly:
\begin{equation}\label{eq:Delta-derivative}
 \Delta_r=r\mathcal W,\qquad \mathcal W=2V+rV_r.
\end{equation}
The boundary coefficients in \eqref{eq:V-division} give
\begin{equation}\label{eq:W-boundary}
 \mathcal W_m(0,0,p)=\frac4{A(p)},\qquad
 \mathcal W_r(0,0,p)=3B_\Delta(0,p)
   =\frac{16c_0(p)}{R'_{p,0,0}(0)A(p)}\ne0.
\end{equation}
Joint smoothness and compactness in \(p\) give a common smaller box
on which \(\mathcal W_r\) has a uniform nonzero margin with sign
\(c_0/A\).

\begin{proposition}[At most one critical point of the separation graph]\label{prop:one-fold-graph}
For fixed \((m,p)\), the graph \(r\mapsto\Delta(r,m,p)\) has at most
one positive critical point on a common interval. Any such point
is nondegenerate. In the outer region \(|\delta|/r^2\le\eta\),
\begin{equation}\label{eq:curvature-factor}
 P_{m,\delta}''(r^2,p)
 =r^{-3}\mathcal U(r,m,\delta/r^2,p)
                  [\delta-\Delta(r,m,p)],
\end{equation}
where \(\mathcal U\) is jointly smooth and nonzero, of sign \(-A\).
For fixed parameters the zeros of \(P''\) in the outer region have total
multiplicity at most two.
\end{proposition}

\begin{proof}
For fixed \((m,p)\), \(\mathcal W\) is strictly monotone and has
at most one positive zero. At that zero,
\(\Delta_{rr}=r\mathcal W_r\ne0\). Hence \(\Delta\) has at most
two monotone branches.

For the factorization, define
\[
 \mathcal U(r,m,v,p)
  =\int_0^1J_v(r,m,V+t(v-V),p)\,dt .
\]
The derivative bound from Lemma~\ref{lem:curvature-profile}
gives its sign and nonvanishing. The exact mean-value identity is
\(J=(v-V)\mathcal U\). Substituting
\(v=\delta/r^2\) in \eqref{eq:J-definition} gives
\eqref{eq:curvature-factor}.
A horizontal line meets this graph in at most two transverse
points, or at its sole positive critical point. Transverse intersections are simple;
an intersection at the nondegenerate critical point has order two.
Changing from \(r>0\) to \(q=r^2\) is a smooth local diffeomorphism,
and the prefactor in \eqref{eq:curvature-factor} never vanishes.
This proves the multiplicity assertion as well as the distinct count.
\end{proof}

We now treat the complement of the outer region, where an active input may
be arbitrarily close to its threshold.

\subsection{The count on the whole interval}\label{sec:inner}

Use the thresholds $0$ and $a_\delta=\Gamma_1^{-1}(\delta)$
from \eqref{eq:passage:true-thresholds}, with the active intervals determined
there. All constants are uniform on the fixed parameter reduction.

Where both passages are active in the threshold layer, compare the difference of the singular
curvatures with their sum. Separation keeps the active inputs at a fixed
relative distance; the difference then dominates both the relative error
and the bounded closing curvature.

Put \(\vartheta=|\delta|\), \(\sigma=\sgn\delta\), and let \(L\ge2\) be fixed.
The strengths satisfy, uniformly,
\[
 A_1(p,m,\delta)=A(p)+O(|m|+\vartheta),\qquad
 A_2(p,m,\delta)=-A(p)+O(|m|+\vartheta).
\]
The smooth root profiles in \eqref{eq:passage:quadratic-jet} yield
\begin{align}
 \Gamma_\nu'(x)
  &=1+\tfrac32 A_\nu\sqrt{x}+O(x),\label{eq:inner-derivative}\\
 \Gamma_\nu''(x)
  &=\tfrac34 A_\nu x^{-1/2}+O(1),\qquad x>0.\label{eq:inner-curvature}
\end{align}
Differentiating the smooth root profiles gives these remainders and their
required derivative bounds. Each $\Gamma_\nu'$ is uniformly positive and
bounded, with active curvature of the sign of its strength.

For \(0<q\le L\vartheta\),
\begin{equation}\label{eq:first-correction}
 |\Gamma_1(q)-q|\le Kq^{3/2}\le K L^{3/2}\vartheta^{3/2}.
\end{equation}
After reducing \(\vartheta\), this is at most \(\vartheta/2\).
The inverse derivative estimates give \(\vartheta/2<a_\vartheta<2\vartheta\).
Every positive penetration in this layer is at most \((L+2)\vartheta\).

\begin{lemma}[Curvature signs in the threshold layer]\label{lem:event-layer}
For every fixed \(L\ge2\), one may choose common sufficiently small
bounds on \(|m|\) and \(\vartheta=|\delta|\) such that, for \(0<|\delta|\),
the signs of \(P''\) are
\begin{equation}\label{eq:event-signs}
 \sgn P''=
 \begin{cases}
 \sgn c,&q<\min(0,a_\delta),\\
 \sgn(\delta A),&\min(0,a_\delta)<q<\max(0,a_\delta),\\
 \sgn(-\delta A),&
       q>\max(0,a_\delta),\quad q\le L|\delta|.
 \end{cases}
\end{equation}
Both passages are active on the last interval. There is no zero of \(P''\) in any
of these open regions. In the active regions
\(|Q''|\ge\kappa_L \vartheta^{-1/2}\) for a uniform \(\kappa_L>0\).
The estimates are uniform up to the thresholds.
\end{lemma}

\begin{proof}
Below both thresholds both passages are inactive, so \(Q=q\) and
\(P''=R''\), with sign \(c\).

If \(\delta=\vartheta>0\), the one-active interval is \(0<q<a_\vartheta\), where
\(Q=\Gamma_1(q)\). Its curvature has sign \(A\), and
\eqref{eq:inner-curvature} gives
\(|Q''|\ge\kappa_L(2\vartheta)^{-1/2}\) after shrinking.
If \(\delta=-\vartheta<0\), that interval is \(-\vartheta<q<0\), where
\(Q=-\vartheta+\Gamma_2(q+\vartheta)\). Its curvature has sign \(-A\) and the same
uniform lower bound.

Where both passages are active, put
\[
 y=\Gamma_1(q)-\delta>0.
\]
The exact chain rule is
\[
 Q''=\Gamma_2'(y)\Gamma_1''(q)
          +\Gamma_2''(y)\Gamma_1'(q)^2.
\]
Equations~\eqref{eq:inner-derivative}--\eqref{eq:first-correction}
therefore give
\begin{equation}\label{eq:relative-inner-error}
 Q''=\tfrac34 A(q^{-1/2}-y^{-1/2})+E,\qquad
 |E|\le C_L(\sqrt \vartheta+\vartheta+|m|)
                 (q^{-1/2}+y^{-1/2}).
\end{equation}
Each error is measured relative to the sum of the two singular
weights, so the formula remains uniform as either positive argument tends
to zero.

By \eqref{eq:first-correction},
\(\sigma(q-y)\ge \vartheta/2\) and \(q,y\le(L+2)\vartheta\). Hence
\[
 -\sigma\frac{q^{-1/2}-y^{-1/2}}{q^{-1/2}+y^{-1/2}}
 =\frac{|q-y|}{(\sqrt q+\sqrt y)^2}
 \ge\frac1{8(L+2)}=:d_L>0.
\]
First reduce \(|m|\), and then \(\vartheta\), so the error in
\eqref{eq:relative-inner-error} is less than half the principal bound.
Thus \(\sgn Q''=\sgn(-\sigma A)\) and
\(|Q''|\ge\kappa_L\vartheta^{-1/2}\), throughout the open two-active layer.

Finally the exact closing-leg chain rule is
\[
 P''=R''(Q)(Q')^2+R'(Q)Q''.
\]
The first term is uniformly bounded, whereas \(R'\) has a positive
margin and the preceding active curvature bound diverges as
\(\vartheta\downarrow0\). A further fixed reduction proves
\eqref{eq:event-signs}.
\end{proof}

\begin{proposition}[Upper bound on the whole interval]\label{prop:common-count}
On a sufficiently small interval $I=(-q_0,q_0)$, common to the family, every
return map $P_\lambda$, $\lambda\in U$, has at most six fixed points.
\end{proposition}

\begin{proof}
First choose the outer box and its fixed \(\eta\) in
Lemma~\ref{lem:curvature-profile}. Set \(L=1/\eta\ge2\), shrinking
\(\eta\) if needed. Apply Lemma~\ref{lem:event-layer} with this one
\(L\): reduce the common strength bound and then the separation bound.
All further reductions are made inside the first-return domain of
Definition~\ref{def:admissible-family}. Choose \(q_0,\delta_0\) so both thresholds remain
strictly inside \(I=(-q_0,q_0)\), for example using
\(\delta_0<q_0/4\) and the uniform inverse derivative bound.

For \(\delta\ne0\), \(P''\) has no zeros in the three regions
of~\eqref{eq:event-signs}. Any remaining zero of \(P''\) has
\(q>L|\delta|>0\), so it lies in the outer region \(|\delta|/q<\eta\).
Proposition~\ref{prop:one-fold-graph} exhausts this region and gives at most
two zeros of \(P''\). At each threshold \(P''\) has a singularity with
nonzero coefficient, so a threshold is not a zero of \(P''\).

The derivative \(P'\) is continuous at the thresholds.
It is strictly monotone on at most five consecutive pieces: the two
thresholds and at most two zeros of \(P''\) are the only possible
divisions. If the signs of \(P''\) on the two sides of a threshold agree,
monotonicity extends through the threshold; if a zero of \(P''\) does not
change sign, it adds no new piece. There is no interval on which \(P''\)
vanishes. Thus \(P'-1\) has at most five zeros. Rolle's theorem for the
\(C^1\) displacement \(D=P-q\), including at the thresholds, gives at most
six zeros.

The case \(\delta=0\) follows from the next corollary.
\end{proof}

\begin{corollary}[The four-point slices]\label{cor:four-slices}
On the same interval the bound on the number of fixed points is four when
\(\delta=0\) or \(m=0\), and two when \(m=\delta=0\).
\end{corollary}

\begin{proof}
For \(\delta=0\), a positive zero of \(P''\) has \(V(r,m,p)=0\).
Equations~\eqref{eq:B-boundary}--\eqref{eq:V-division} give
\[
 V_r(0,0,p)=B_\Delta(0,p)=\frac{16c_0(p)}{3k(p)A(p)}\ne0.
\]
Compactness in \(p\) gives a common nonzero margin for \(V_r\)
on a smaller box. Thus, for fixed \(m,p\), there is at most one
positive zero of \(V\), and it is simple. Formula~\eqref{eq:curvature-factor}
and the local diffeomorphism \(q=r^2\), \(r>0\), preserve its order. Together with at most one sign change at the aligned threshold,
this gives at most three monotonicity pieces for \(P'\) and four fixed points.

For \(m=0,\delta\ne0\), the separation graph is
\(\Delta=r^3B_\Delta(r,p)\), with sign \(c_0/A\) and nonzero derivative.
It has at most one curvature intersection, necessarily simple, and
one is possible only if \(c_0\delta A>0\).
When this inequality holds, the first two signs in
\eqref{eq:event-signs} agree, leaving one threshold change and
at most one outer change. If \(c_0\delta A<0\), there is no outer
intersection and at most two threshold changes. Again \(P'\)
has at most three monotonicity pieces.

At \(m=\delta=0\), equation~\eqref{eq:center-curvature} and the
inactive closing map give the same strict sign \(c_0\) throughout
the interval. The displacement is strictly convex or strictly
concave, including across \(q=0\), and has at most two zeros.
\end{proof}

\begin{corollary}[Open six-cycle regions]\label{cor:six-open}
In every admissible family, the parameters \(\lambda\in U\) for which
\(P_\lambda\) has six simple fixed points in \(I\) form an open set.
\end{corollary}

\begin{proof}
Choose six disjoint compact intervals isolating the simple zeros of \(D\).
Joint \(C^1\) dependence, including at the thresholds, preserves the opposite
endpoint signs and the nonzero derivative signs for nearby parameters, so each
interval keeps one simple zero. Proposition~\ref{prop:common-count} excludes
further zeros in \(I\).
\end{proof}

Propositions~\ref{prop:single-scale-six} and~\ref{prop:sharp-slice-attainment}
below show that all four bounds are attained under~\eqref{eq:rank}.


\subsection{The cubic critical value}

When the separation graph has a positive critical point, its height determines the
two-intersection window. We compute this height and its cubic separation
scale for the saturation argument.

\begin{proposition}[The critical value of the separation graph]
\label{prop:cubic-curvature-window}
After a common further reduction, the signed critical point of the separation
graph is
\[
 r_{\rm c}(m,p)=-\frac{k(p)m}{4c_0(p)}+O(m^2),\qquad r_{\rm c}(0,p)=0.
\]
It is smooth through zero. A positive critical point exists exactly
when \(mc_0(p)<0\), and then it equals \(r_{\rm c}(m,p)\).
Its critical value has a smooth division
\begin{equation}\label{eq:cubic-critical-value}
 \delta_{\rm c}(m,p)=\Delta(r_{\rm c}(m,p),m,p)=m^3\kappa_{\rm c}(m,p),
 \qquad \kappa_{\rm c}(0,p)=\frac{k(p)^2}{24A(p)c_0(p)^2}.
\end{equation}
The remainder and all reductions are uniform in \(p\).
On a fixed interval \(0<r<\sqrt{q_0}\), reduced as below, \(P''\) has
exactly two simple zeros in the outer region if and only if
\begin{equation}\label{eq:curvature-window}
 mc_0(p)<0,\qquad 0<\delta/\delta_{\rm c}(m,p)<1.
\end{equation}
At \(\delta=\delta_{\rm c}\), \(mc_0<0\), it instead has one double zero
in the outer region.
\end{proposition}

\begin{proof}
Apply the implicit function theorem to \(\mathcal W=0\).
The boundary derivatives \eqref{eq:W-boundary} and the uniform margin
give a unique smooth signed solution \(r_{\rm c}=m\gamma\), with
\(\gamma(0,p)=-k/(4c_0)\). Its positive part exists precisely for
\(mc_0<0\).
Substitution into \(\Delta=r^2(mH_\Delta+rB_\Delta)\) gives the smooth quotient
\[
 \kappa_{\rm c}=\gamma^2\{H_\Delta(m\gamma,m,p)+\gamma B_\Delta(m\gamma,p)\}.
\]
The boundary coefficients yield \(\gamma B_\Delta=-4/(3A)\), hence
\eqref{eq:cubic-critical-value}. Thus \(\kappa_{\rm c}\) has the sign of
\(A\); when a positive critical point exists, the critical value has sign
\(-c_0/A\).

First fix \(r_0=\sqrt{q_0}\) inside the common root box. Then reduce
\(|m|\), uniformly in \(p\), so \(0<r_{\rm c}<r_0/2\) when
\(mc_0<0\), and
\[
 |mH_\Delta(r_0,m,p)|<\tfrac12r_0|B_\Delta(r_0,p)|.
\]
Thus \(\Delta(r_0,m,p)\) has sign \(c_0/A\), opposite to the
nonzero critical value. Also require \(|\delta_{\rm c}|\) to lie
within the common separation bound. The unique positive critical point and
\(\Delta(0)=0\) give exactly two transverse intersections with the
line \(\delta\) precisely for \eqref{eq:curvature-window}.
At each intersection \(|\delta|/r^2=|V(r,m,p)|<\eta\), so both
lie in the outer region. Conversely two distinct outer
intersections force the unique positive critical point between them and place
\(\delta\) strictly between zero and its critical value.
The factorization \eqref{eq:curvature-factor} transfers their orders
to \(P''\). At the critical value, the nondegenerate critical point gives
\(D''=D'''=0\) and \(D''''\ne0\).
\end{proof}

On the curve \(\delta=\delta_{\rm c}\) the conditions \(D=D'=0\) are not
imposed, so this curve is not a fold curve of periodic orbits.

Figure~\ref{fig:curvature-graph} shows a separation graph with one
positive critical point.
% Schematic drawing of one possible shape of the separation graph.
\begin{figure}[!htbp]
  \centering
  \begin{tikzpicture}[x=1cm,y=1cm,font=\normalsize,
      axis/.style={line width=.55pt,-{Stealth[length=2mm]}},
      graph/.style={line width=.9pt},
      guide/.style={line width=.45pt,densely dashed},
      projection/.style={line width=.4pt,densely dotted}]
    \node at (5.1,3.0)
      {$\Delta(r,m,p)=r^2V(r,m,p)=r^2\bigl[mH_\Delta(r,m,p)+rB_\Delta(r,p)\bigr]$};

    \draw[axis] (0,-1.15) -- (0,2.55)
      node[above] {$\delta$};
    \node[anchor=east,align=right] at (-.22,2.12)
      {threshold\\separation};
    \draw[axis] (0,0) -- (10.05,0)
      node[below] {$r=\sqrt q$};
    \node[below left] at (0,0) {$0$};

    % The two selected intersections lie on this hand-drawn guide.
    \draw[graph]
      (0,0)
      .. controls (.8,0) and (1.55,.45) .. (2.65,1.1)
      .. controls (3.75,1.75) and (4.55,2.05) .. (5.35,2.05)
      .. controls (6.05,2.05) and (7.1,1.72) .. (7.8,1.1)
      .. controls (8.45,.52) and (9.0,-.32) .. (9.32,-.94);
    \node[anchor=east] at (8.5,-.94) {$\Delta(r,m,p)$};

    \draw[guide] (0,1.1) -- (9.5,1.1)
      node[anchor=west] {$\delta=\text{constant}$};
    \draw[projection] (2.65,0) -- (2.65,1.1);
    \draw[projection] (7.8,0) -- (7.8,1.1);
    \draw[projection] (5.35,0) -- (5.35,2.05);
    \fill (2.65,1.1) circle[radius=2pt];
    \fill (7.8,1.1) circle[radius=2pt];
    \node[below] at (2.65,0) {$r_1$};
    \node[below] at (7.8,0) {$r_2$};
    \node[below] at (5.35,0) {$r_{\rm c}$};
    \node[above] at (5.35,2.08) {positive critical point};
  \end{tikzpicture}
  \caption{The separation graph $\delta=\Delta(r,m,p)$ for fixed $(m,p)$ with
    $c_0<0$, $A>0$ and $m>0$ (schematic). The graph has at most one positive
    critical point, which is nondegenerate
    (Proposition~\ref{prop:one-fold-graph}). Its intersections with a
    horizontal line are the zeros of $P''$ in the outer region $q>L|\delta|$,
    by~\eqref{eq:curvature-factor}; for other parameters the graph can be
    monotone. The resulting bound on fixed points is
    Proposition~\ref{prop:common-count}.}
  \label{fig:curvature-graph}
\end{figure}



\subsection{Attainment of the bounds}\label{subsec:single-scale-sharpness}

To attain six fixed points we choose a quartic with four simple positive
roots; its coefficients determine a quadratic with two simple negative roots.
We continue the two groups of roots in disjoint state intervals and then
treat the slices \(m=0\) and \(\delta=0\).

Assume~\eqref{eq:rank} and use the standard unfolding
\(\lambda=\Lambda(a,b,m,\delta)\) of Section~\ref{sec:introduction}, with
\(\lambda_0=\Lambda(0,0,0,0)\). Write
\[
 A_0=A_1(\lambda_0),\quad
 c_0=\tfrac12R_{\lambda_0}''(0),\quad
 B=|A_0|,\quad C=|c_0|,
\]
\[
 s=-\sgn c_0,\qquad \tau=\sgn(sA_0).
\]
We first allow an arbitrary rescaled strength sum, including zero:
\begin{equation}\label{eq:weighted-general-controls}
 \lambda_\eps=\Lambda(\bar a\eps^4,\bar b\eps^2,
                         \bar m\eps,\bar d\eps^3),\qquad
 \theta=(\bar a,\bar b,\bar m,\bar d),\qquad \eps>0.
\end{equation}

\begin{lemma}[Rescaled displacement]\label{lem:weighted-quartic}
Let \(\Omega_+\Subset(0,\infty)\), \(\Omega_-\Subset(-\infty,0)\), and let
\(\theta\) range over a fixed compact set. For sufficiently small
\(\eps>0\), the rescaled displacements satisfy, for \(z\in\Omega_+\)
and \(x\in\Omega_-\),
\begin{align}
 \mathcal F_\eps(z)&:=\eps^{-4}D_{\lambda_\eps}(\eps^2z^2)
 =F_\theta(z)+\eps E_+(\eps,z,\theta),\label{eq:quartic-limit}\\
 F_\theta(z)&=\bar a+\bar bz^2+\bar mz^3+c_0z^4+\tfrac32A_0\bar dz,\notag\\
 \mathcal G_\eps(x)&:=\eps^{-4}D_{\lambda_\eps}(\eps^2x)
 =G_\theta(x)+\eps E_-(\eps,x,\theta),
 \label{eq:quartic-negative-limit}\\
 G_\theta(x)&=\bar a+\bar bx+c_0x^2.\notag
\end{align}
The remainders extend jointly smoothly to signed \(\eps\) near zero.
For every prescribed finite integer \(n\), uniformly on these compacts,
\begin{equation}\label{eq:quartic-mixed-norm}
 \max_{j+|\gamma|\le n}
 \left\|\partial_z^j\partial_\theta^\gamma
                 (\mathcal F_\eps-F_\theta)\right\|_\infty\le C_n\eps,
\end{equation}
and similarly with \(x,\mathcal G_\eps\). The compact parameter set may
contain points with \(\bar m=0\).
\end{lemma}

\begin{proof}
Use the exact identities of Lemma~\ref{lem:exact-value-division}.
Substitute \(r=\eps z\), \(v=\bar d\eps/z^2\) into
\eqref{eq:quartic-Hadamard}; the profiles are evaluated at
\(\delta=r^2v=\bar d\eps^3\). Then
\[
 Q-q=\eps^4\Xi,\qquad
 \Xi=\bar mz^3+\bar dzH+\eps z^5L_0+\tfrac34\bar m^2\eps^2z^4.
\]
With \(Y=z^2+\eps^2\Xi\), the closing expansion gives
\begin{equation}\label{eq:quartic-exact-residual}
 \mathcal F_\eps=
 \bar a+\bar bY+\Xi+c(\lambda_\eps)Y^2
                   +\eps^2Y^3S(\eps^2Y,\lambda_\eps).
\end{equation}
Since \(z\) stays away from zero, all arguments remain in the common
root box, whose positive square-root factor stays bounded away from zero.
At \(\eps=0\), \(H=3A_0/2\) and \(c=c_0\), giving \(F_\theta\).
Taylor integral division in \(\eps\) proves the joint remainder and
every finite mixed estimate in \eqref{eq:quartic-mixed-norm}, without
dividing by \(\bar m\). State derivatives are taken at fixed
\((\eps,\theta)\), hence at fixed \(\lambda_\eps\); the variable
\(v=\bar d\eps/z^2\) keeps the separation \(r^2v=\bar d\eps^3\) fixed.

For negative \(x\) in a fixed compact set, \(q=\eps^2x\) is below
both thresholds for small positive \(\eps\), since \(\delta=O(\eps^3)\) and
\(|a_\delta|\le2|\delta|\). Thus \(Q=q\) exactly, and
\eqref{eq:quartic-closing-division} gives
\[
 \mathcal G_\eps=\bar a+\bar b x+c(\lambda_\eps)x^2
                   +\eps^2x^3S(\eps^2x,\lambda_\eps).
\]
This proves the negative formula and its joint smooth extension.
\end{proof}

For the six-point construction specialize to
\(\bar a=s\alpha_4\), \(\bar b=s\bsc\), \(\bar m=s\) and
\(\bar d=\tau\dsc\), giving
\begin{equation}\label{eq:quartic-weights}
 m=s\eps,\qquad \delta=\tau \dsc\eps^3,\qquad
 b=s\bsc\eps^2,\qquad a=s\alpha_4\eps^4,\qquad \eps>0.
\end{equation}
Here \(\alpha_4,\bsc,\dsc\) are rescaled parameters.
Multiplying the rescaled displacements by \(s\), and using
\(s\tau A_0=B\), \(sc_0=-C\), gives
\(\mathcal F_{4,\eps}:=s\mathcal F_\eps=F_0+O(\eps)\) and
\(\mathcal G_{4,\eps}:=s\mathcal G_\eps=G_0+O(\eps)\), where
\[
 \begin{aligned}
 F_0(z)&=\alpha_4+\bsc z^2+z^3-Cz^4+\tfrac32B\dsc z,\\
 G_0(x)&=\alpha_4+\bsc x-Cx^2.
 \end{aligned}
\]
These expansions have all the mixed estimates of
Lemma~\ref{lem:weighted-quartic}.

Take distinct positive numbers \(r_i\), \(1\le i\le4\), with
\(\sum r_i=1/C\), and let \(e_j\) be their elementary symmetric
functions. Choosing
\begin{equation}\label{eq:quartic-Vieta}
 \alpha_4=-Ce_4,\qquad \bsc=-Ce_2,\qquad \dsc=\frac{2C}{3B}e_3
\end{equation}
gives simultaneously
\[
 F_0(z)=-C\prod_{i=1}^4(z-r_i),\qquad
 G_0(x)=-C(x^2+e_2x+e_4).
\]
The six pair products have product \(e_4^3\); AM--GM yields
\(e_2\ge6\sqrt{e_4}\). Thus \(e_2^2-4e_4>0\), and the quadratic
has two distinct negative roots. Hence the choice~\eqref{eq:quartic-Vieta}
gives four simple positive roots of \(F_0\) and two simple negative roots
of \(G_0\).

\begin{proposition}[Six fixed points]
\label{prop:single-scale-six}
For every admissible family satisfying \eqref{eq:rank}, the standard
unfolding has six simple fixed points of order $\eps^2$ for all parameters
\eqref{eq:quartic-weights} with rescaled parameters in an open set and all
sufficiently small $\eps>0$. Their stability alternates: three are attracting
and three repelling.
\end{proposition}

\begin{proof}
Choose any distinct positive $r_i$ with sum $1/C$ and the parameters
\eqref{eq:quartic-Vieta}. The four roots of $F_0$ and the two negative
roots of $G_0$ have strict derivative margins. Lemma~\ref{lem:weighted-quartic}
continues them on six disjoint isolating intervals, uniformly over a small
open set of rescaled parameters. The same joint estimates allow \(\bar m\) to
vary near \(s\), so this set may be taken open in all four rescaled
parameters. Fixed rescaled parameters give a smooth parameter path tending to
the reference, and Proposition~\ref{prop:common-count} excludes additional
fixed points. At these parameters, for negative states $D_q=s\eps^2G_0'(x)+O(\eps^3)$, and for positive
states $D_q=s\eps^2F_0'(z)/(2z)+O(\eps^3)$. In state order these signs
are $s(+,-,+,-,+,-)$. Since the return is increasing and close to the
identity, they give three attracting and three repelling hyperbolic cycles;
these signs persist on the open set of rescaled parameters.
\end{proof}

\begin{proposition}[Attainment on the slices]
\label{prop:sharp-slice-attainment}
For every admissible family satisfying \eqref{eq:rank}, the standard unfolding
has four simple fixed points on \(\delta=0\), four on \(m=0\), and two on
\(m=\delta=0\). Each configuration occurs in \(I\) for every sufficiently
small \(\eps>0\) along a smooth parameter path tending to the reference, and
persists on an open set of rescaled parameters within the corresponding slice.
\end{proposition}

\begin{proof}
Use the weights \eqref{eq:weighted-general-controls}. The following table
specifies the paths and the limiting simple roots; positive roots use the
coordinate $z$, and negative roots the coordinate $x$:
\begin{equation}\label{eq:sharp-slice-controls}
 \begin{array}{c|c|c|c}
 \text{slice}&(\bar a,\bar b,\bar m,\bar d)/(sC)&z>0&x<0\\ \hline
 \delta=0&(36/11,-85/11,60/11,0)&1,2,3&x_\delta\\
 m=0&(36,25,0,-40/A_0)&1,2,3&x_m\\
 m=\delta=0&(1,0,0,0)&1&-1
 \end{array}
\end{equation}
Here $x_\delta=(-85-\sqrt{8809})/22$ and
$x_m=(25-\sqrt{769})/2$. For the first two rows, direct substitution in
Lemma~\ref{lem:weighted-quartic} gives respectively
\[
 \begin{aligned}
 F_\theta/(sC)&=-(z-1)(z-2)(z-3)(z+6/11),&
 G_\theta/(sC)&=36/11-85x/11-x^2,\\
 F_\theta/(sC)&=-(z-1)(z-2)(z-3)(z+6),&
 G_\theta/(sC)&=36+25x-x^2.
 \end{aligned}
\]
Only the positive roots of $F_\theta$ correspond to fixed points; the
negative ones come from $G_\theta$. The estimates of
Lemma~\ref{lem:weighted-quartic} include $\bar m=0$.

For the third row, \eqref{eq:center-displacement} and the root-profile
remainder estimates give, uniformly on compact $x$-intervals including zero,
\[
 \frac{D_{\lambda_\eps}(\eps^2x)}{sC\eps^4}
 =1-x^2+O_{C^1}(\eps).
\]

Continue the listed simple roots by Lemma~\ref{lem:weighted-quartic} and
the isolating argument in the proof of Proposition~\ref{prop:single-scale-six}
on small relatively open sets of rescaled parameters with compact closures:
$\bar d=0$ in the first row, $\bar m=0$ in the second, and $\bar m=\bar d=0$
with $(\bar a,\bar b)$ near $(sC,0)$ in the third. The rescaled thresholds
are $0$ and $a_\delta/\eps^2$, with
$|a_\delta|/\eps^2\le2|\bar d|\eps$. One reduction of
$\eps$ places the positive intervals where both passages are active, the
negative ones below both thresholds, the parameters in the domain of
$\Lambda$, and the rescaled intervals in $I$.
Corollary~\ref{cor:four-slices} excludes other roots in $I$, including at the
thresholds, and $\Lambda$ gives smooth parameter paths and relative
persistence.
\end{proof}

The positive quartet in the six-cycle construction is the starting point for
Section~\ref{sec:main-saturation}.
